% !TEX root = ../algebra_02.tex

\section{Теорема о факторгруппе факторгруппы}

\begin{Theorem}{О факторгруппе факторгруппы}{}
	Пусть $H, P \triangleleft G$ и $H \subset P$. Тогда:
	\begin{enumerate}
		\item $H \triangleleft P$;
		\item $P/H \triangleleft G/H$;
		\item $(G/H)/(P/H) \cong G/P$.
	\end{enumerate}
\end{Theorem}

\begin{proof}
	1. Так как $H \triangleleft G$, то для любых $g \in G$ и $h \in H$ выполнено $ghg^{-1} \in H$. Поскольку $P \subset G$, это условие автоматически выполняется и для всех $g \in P$, что означает $H \triangleleft P$.

	2--3. Рассмотрим отображение $\pi_P \colon G \to G/P$, заданное правилом $g \mapsto gP$. Так как $H \subset P = \ker \pi_P$, существует индуцированный гомоморфизм:
	\[
	\widetilde{\pi}_P \colon G/H \to G/P, \quad gH \mapsto gP.
	\]
	Очевидно, что $\Im \widetilde{\pi}_P = G/P$. Найдем ядро этого гомоморфизма:
	\[
	\ker \widetilde{\pi}_P = \{ gH \mid gP = eP \} = \{ gH \mid g \in P \} = P/H.
	\]
	Следовательно, ядро является нормальной подгруппой: $P/H \triangleleft G/H$.
	Применяя теорему о гомоморфизме к $\widetilde{\pi}_P$, получаем:
	\[
	(G/H)/(P/H) \cong G/P.
	\]
\end{proof}

\begin{proof}[Альтернативное доказательство пункта 3 (прямое построение изоморфизма)]
	Построим изоморфизм напрямую. Рассмотрим отображение
	\[
	\varphi \colon G/P \to (G/H)/(P/H), \quad \varphi(gP) = (gH)(P/H).
	\]
	Проверим корректность. Пусть $g_1 P = g_2 P$. Тогда $g_2^{-1} g_1 \in P$. Нам нужно показать, что $(g_1 H)(P/H) = (g_2 H)(P/H)$, что эквивалентно $(g_2 H)^{-1} (g_1 H) \in P/H$. 
	
	Заметим, что $(g_2 H)^{-1} (g_1 H) = (g_2^{-1} g_1) H$. Так как $g_2^{-1} g_1 \in P$, смежный класс $(g_2^{-1} g_1) H$ принадлежит $P/H$. Следовательно, отображение корректно определено.

	Проверим гомоморфность:
	\[
	\varphi((g_1 P)(g_2 P)) = \varphi(g_1 g_2 P) = (g_1 g_2 H)(P/H) = ((g_1 H)(P/H))((g_2 H)(P/H)) = \varphi(g_1 P)\varphi(g_2 P).
	\]

	Проверим сюръективность. Произвольный элемент из $(G/H)/(P/H)$ имеет вид $(gH)(P/H)$ для некоторого $g \in G$. Это в точности равно $\varphi(gP)$, откуда следует сюръективность.

	Проверим инъективность. Пусть $gP \in \ker \varphi$. Тогда $\varphi(gP) = P/H$, что означает $(gH)(P/H) = P/H$, то есть $gH \in P/H$. Значит, существует такой элемент $p \in P$, что $gH = pH$, откуда $p^{-1}g \in H$. Поскольку $H \subset P$ и $p \in P$, имеем $p(p^{-1}g) \in P$, то есть $g \in P$. Следовательно, $gP = P$, ядро тривиально.

	Отображение $\varphi$ является биективным гомоморфизмом, то есть изоморфизмом.
\end{proof}
